% questions.tex —— 1980 年·全国卷（理科） · 题目内容片段（无答案）
% 由 tools/generate_projects.py 依来源总账自 1980/1980_数学真题_理_试题.md 生成，勿手改（改源 MD 后重跑）


\gkpart{理科数学}

\begin{gktip}
\textbf{来源说明}：本题面依据固定原卷 PDF《1980 理》第 1 页（印刷页码 15）逐题转录。源卷为单页题面，未附答案页；配套答案与解析由本库独立推导，\textbf{非官方标准答案}。原卷为三栏排印：第 1—4 题在第一栏，第 5—7 题在第二栏，第 8—9 题与附加题在第三栏，本文按题号顺序转录。正式题为第 1—9 题；卷末「附加题」栏一题在原卷中沿用大题号印作「10．」，本库作为附加题 10 保留在第 9 题之后，不计入正式题量。第 5 题附立体示意图，按源图重绘为 TikZ 矢量源。第 6、8 两题原卷印刷与本库核对结论不一致，见各题源注。
\end{gktip}

\begin{question}
\gkti{1}将多项式 $x^5y-9xy^5$ 分别在下列范围内分解因式：

（1） 有理数范围；

（2） 实数范围；

（3） 复数范围。

\end{question}
\begin{question}
\gkti{2}半径为 $1$、$2$、$3$ 的三个圆两两外切。证明：以这三个圆的圆心为顶点的三角形是直角三角形。

\end{question}
\begin{question}
\gkti{3}用解析几何方法证明：三角形的三条高线交于一点。

\end{question}
\begin{question}
\gkti{4}证明对数换底公式：$\log_b N=\dfrac{\log_a N}{\log_a b}$（$a$、$b$、$N$ 都是正数，$a\neq1$，$b\neq1$）。

\Needspace{0.4\textheight}
\end{question}
\begin{question}
\gkti{5}直升飞机上一点 $P$ 在地面 $M$ 上的正射影是 $A$，从 $P$ 看地面上一物体 $B$（不同于 $A$）。直线 $PB$ 垂直于飞机窗玻璃所在的平面 $N$（如图）。证明：平面 $N$ 必与平面 $M$ 相交，且交线 $l$ 垂直于 $AB$。
\nopagebreak[4]
\begin{center}% 1980 年全国卷（理科）第 5 题：直升飞机上点 P、地面 M 与飞机窗玻璃平面 N 的立体示意图
% 按源卷第 1 页第二栏源图布局重绘：地面 M 为透视平行四边形；斜置平面 N 与 M 相交，交线 l 画在 M 上；
% P 在 M 上的正射影为 A（竖直虚线 PA）；B 为地面上另一点，连 AB；被不透明平面 N 遮住的部分画虚线
\begin{tikzpicture}[scale=1.0,line join=round]
  % 地面 M（透视平行四边形）
  \coordinate (M1) at (2.0,2.9);
  \coordinate (M2) at (13.5,2.9);
  \coordinate (M3) at (16.1,6.8);
  \coordinate (M4) at (4.6,6.8);
  % 平面 N（斜置平行四边形：Q1 上、Q2 左、Q3 下、Q4 右）
  \coordinate (Q1) at (8.4,10.3);
  \coordinate (Q2) at (5.4,7.7);
  \coordinate (Q3) at (10.9,3.9);
  \coordinate (Q4) at (13.9,6.5);
  % 关键点
  \coordinate (P) at (9.1,11.6);
  \coordinate (A) at (9.1,4.9);
  \coordinate (B) at (4.4,4.9);
  % 交线 l 与 M 的后棱、N 的边界相交处
  \coordinate (L0) at (7.546,6.8);
  \coordinate (L1) at (12.057,2.9);
  \coordinate (K1) at (6.70,6.8);
  \coordinate (K2) at (13.47,6.8);
  \coordinate (V1) at (9.1,9.816);
  \coordinate (V2) at (9.1,5.144);
  \coordinate (H1) at (9.452,4.9);
  % 地面 M
  \draw[thick] (M1)--(M2)--(M3);
  \draw[thick] (M4)--(M1);
  % M 的后棱：被平面 N 遮住的一段画虚线
  \draw[thick] (M4)--(K1);
  \draw[dashed,thick] (K1)--(K2);
  \draw[thick] (K2)--(M3);
  % 平面 N
  \draw[thick] (Q1)--(Q2)--(Q3)--(Q4)--cycle;
  % 交线 l（M 与 N 的交线）
  \draw[thick] (L0)--(L1);
  % 直线 AB：被 N 遮住的一段画虚线
  \draw[thick] (B)--(H1);
  \draw[dashed,thick] (H1)--(10.8,4.9);
  % P 到 M 的垂线 PA（中段为 N 所遮）
  \draw[thick] (P)--(V1);
  \draw[dashed,thick] (V1)--(V2);
  \draw[thick] (V2)--(A);
  % 点与标注
  \fill (P) circle (1.3pt) node[right=2pt] {$P$};
  \fill (A) circle (1.3pt) node[below=2pt] {$A$};
  \fill (B) circle (1.3pt) node[below left=1pt] {$B$};
  \node at (9.35,8.15) {$N$};
  \node at (14.75,6.05) {$M$};
  \node at (12.35,2.55) {$l$};
\end{tikzpicture}
\end{center}
\nopagebreak[4]
\end{question}
\begin{question}
\gkti{6}设三角函数 $f(x)=\sin\left(\dfrac{k\pi}{5}+\dfrac{\pi}{3}\right)$，其中 $k\neq0$。

\begin{gktip}
\textbf{源注}：原卷（含 PDF 文本层）此处印作 $\sin\left(\frac{k\pi}{5}+\frac{\pi}{3}\right)$，右端与 $x$ 无关；按字面理解 $f(x)$ 是常函数，无最小正周期，第（2）问所说的「随自变量 $x$ 变化」也无从成立。通行题解此处作 $f(x)=\sin\left(\frac{kx}{5}+\frac{\pi}{3}\right)$（角频率为 $k/5$）。本题答案按通行版本给出，推导与两种口径的差别见解析正文。
\end{gktip}

（1） 写出 $f(x)$ 极大值 $M$、极小值 $m$ 与最小正周期 $T$；

（2） 试求最小的正整数 $k$，使得当自变量 $x$ 在任意两个整数间（包括整数本身）变化时，函数 $f(x)$ 至少有一个值是 $M$ 与一个值是 $m$。

\end{question}
\begin{question}
\gkti{7}$CD$ 为直角 $\triangle ABC$ 中斜边 $AB$ 上的高，已知 $\triangle ACD$、$\triangle CBD$、$\triangle ABC$ 的面积成等比数列，求 $\angle B$（用反三角函数表示）。

\end{question}
\begin{question}
\gkti{8}已知 $0<\alpha<\pi$，证明：$2\sin\alpha\leqslant\cot\dfrac{\alpha}{2}$；并讨论 $\alpha$ 为何值时等号成立。

\begin{gktip}
\textbf{源注}：原卷此处印作 $2\sin\alpha$。该式在 $\alpha>\frac{\pi}{3}$ 时不成立（如 $\alpha=\frac{\pi}{2}$ 时左端为 $2$、右端为 $1$），故原卷此处印刷有误；通行题解作 $2\sin 2\alpha\leqslant\cot\frac{\alpha}{2}$，此时不等式对一切 $0<\alpha<\pi$ 成立且等号在 $\alpha=\frac{\pi}{3}$ 取得。本库按通行版本作答，推导见解析正文。
\end{gktip}

\end{question}
\begin{question}
\gkti{9}抛物线的方程是 $y^2=2x$，有一个半径为 $1$ 的圆，圆心在 $x$ 轴上运动。问这个圆运动到什么位置时，圆与抛物线在交点处的切线互相垂直。

（注：设 $P(x_0,y_0)$ 是抛物线 $y^2=2px$ 上一点，则抛物线在 $P$ 点处的切线斜率是 $\dfrac{p}{y_0}$。）

\textbf{附加题 10．} 设直线 $l$ 的参数方程是 $\begin{cases}x=t\\ y=b+mt\end{cases}$（$t$ 是参数），椭圆 $E$ 的参数方程是 $\begin{cases}x=1+a\cos\theta,\ (a\neq0)\\ y=\sin\theta\end{cases}$（$\theta$ 是参数），问 $a$、$b$ 应满足什么条件，使得对于任意 $m$ 值来说，直线 $l$ 与椭圆 $E$ 总有公共点。

\end{question}